\chapter*{Who This Book Is For}

\section*{Who This Book Is For}

This book is for readers who want to understand how enterprise agentic-AI and AI-assisted decision-support systems are designed---not merely how to call a model API or assemble a framework demonstration. Its combination of business context, requirements discovery, pipeline logic, architecture, data modeling, evaluation, and operations is intended for:

\begin{itemize}
\item \textbf{AI and software engineers} building document intelligence, retrieval, agentic workflows, and production decision-support applications.
\item \textbf{Freshers and early-career developers} who know the basics of Python, APIs, databases, or machine learning and want to see how those pieces fit together in industry systems.
\item \textbf{Technical-interview candidates} practising end-to-end system design, requirement discovery, trade-offs, reliability, data modeling, and governance.
\item \textbf{Consultants and solution architects} translating ambiguous financial-services problems into credible solution designs.
\item \textbf{Data scientists and ML engineers} moving from models and notebooks toward deployable, observable, auditable services.
\item \textbf{Fintech and BFSI professionals} who need a practical account of what an LLM can safely contribute to regulated workflows.
\end{itemize}
You do not need to be a credit analyst, compliance officer, or portfolio manager before starting. Chapters 1 and 2 establish the technical and business vocabulary used in the cases. Familiarity with basic programming concepts will help with the low-level design discussions, but the dialogue explains why the components exist before discussing how to implement them.
